
\section{Experiments}
\label{sec:experiments}

We compare \bocas to the following four baselines: 
\gpucb, the \gpeis criterion in \bayo~\citep{jones98expensive}, 
\mfgpucbones from~\citet{kandasamy16mfbo} and
\mfsko: the multi-fidelity sequential kriging optimisation method
from~\citet{huang06mfKriging}.
All methods are based on GPs.
The first two are not multi-fidelity methods, while the last two are finite multi-fidelity
methods\footnote{%
To our knowledge, the only other work that applies to continuous approximations
is~\citet{klein2015towards} which was developed specifically for hyper-parameter tuning.
Further, their implementation is not made available and is not straightforward to
implement.
}.
For the latter,

In all our experiments we use the SE kernel for the GPs.
We have described some implementation details for all these methods in
Appendix~\ref{sec:appImplementation}.


\subsection{Synthetic Experiments}
\label{sec:synthetic}

\insertToyResultsFigure

The results for the first set of synthetic experiments are given in Fig.~\ref{fig:toy}.
The title of each figure states the function used, and the dimensionalities $p,d$ of the
fidelity space and domain.
In all case, the fidelity space was taken to be $\Zcal=[0,1]^p$ with
$\zhf = \one_p = [1,\dots,1]\in\RR^p$
being the most expensive fidelity.
For \mfgpucbones and \mfsko, we used $3$ fidelities ($2$ approximations) where the
approximations were obtained at $z=0.333\one_p$ and $z=0.667\one_p$ points in the fidelity
space.
To reflect the setting in our experiments, we add Gaussian noise to the function value
when observing $g$ at any $(z,x)$.
This makes the problem more challenging than standard global optimisation problems
where function evaluations are not noisy.
The functions $g$, the cost functions $\cost$ and the noise variances $\eta^2$ are given
in Appendix~\ref{sec:appExperiments}.

The first two figures in Fig.~\ref{fig:toy} are simple sanity checks.
In both cases, $\Zcal=[0,1]$ and $\Xcal=[0,1]$ and the functions were sampled from GPs. 
The GP was made known to all methods which used them in picking the next point.
In the first figure, we used an SE kernel for $\phix$ with bandwidth 

In the remaining experiments, we use some standard benchmarks for global optimisation.
We modify them to obtain the function $g$ and add noise to the observations.
As the kernel and other GP hyper-parameters are unknown, we learn them by maximising
the marginal likelihood every $25$ iterations.
This is a common heuristic used in the \bayo{} literature.
We outperform all methods on all problems
except in the case of the Borehole function where \mfgpucbones does better.


\subsection{Maximum Likelihood Estimation with Type Ia Supernovae}
\label{sec:sn}

\subsection{Support Vector Classification with $20$ news groups}
\label{sec:sn}

\insertRealResultsFigure

